\chapter*{Abstract}
La segmentazione automatica dei tumori cerebrali in immagini di risonanza magnetica rappresenta un problema di grande interesse, poiché riduce la necessità di annotare manualmente le lesioni e velocizza le diagnosi. In questo lavoro viene proposta e sviluppata in ambiente MATLAB una pipeline di segmentazione basata su tecniche classiche di \textit{Image Processing}, con l’obiettivo di confrontare l’efficacia di approcci monomodali e multimodali.

La metodologia integra una fase preliminare di analisi esplorativa dei dati, seguita da pre-processing con selezione automatica della slice più significativa, normalizzazione delle intensità, filtraggio mediano del rumore e costruzione di rappresentazioni multimodali mediante fusione pesata delle sequenze \textit{FLAIR}, \textit{T1c} e \textit{T2}. La segmentazione è realizzata attraverso una strategia a cascata che combina sogliatura multilivello di \textit{Otsu}, \textit{Region Growing} con seme automatico e \textit{Marker-Controlled Watershed}, mentre il post-processing morfologico consente il raffinamento delle maschere ottenute.

Le prestazioni sono state valutate mediante le metriche \textit{Dice Score}, \textit{Sensitivity} e \textit{Precision} per singolo paziente e in forma aggregata sull’intero dataset. I risultati mostrano che la configurazione multimodale \textit{Fusion3} risulta la più efficace, con valori medi pari a \textit{Dice Score} \(= 0.806\), \textit{Sensitivity} \(= 0.744\) e \textit{Precision} \(= 0.935\), superando la migliore configurazione monomodale, \textit{FLAIR} e bimodale \textit{Fusion2}. Le sequenze \textit{T2} e \textit{T1c}, considerate singolarmente, mostrano invece prestazioni inferiori, in quanto meno rappresentative della lesione.

Nel complesso, i risultati confermano che l’integrazione di informazioni complementari provenienti da più sequenze MRI può migliorare l’accuratezza della segmentazione rispetto all’uso di singole modalità. Il lavoro mostra inoltre come approcci tradizionali, seppur meno complessi rispetto ai metodi basati su \textit{Deep Learning}, possano comunque fornire risultati significativi nel contesto applicativo analizzato.